\notechapter{N-20220716}{Conic Exercises: Lines, Ellipses, Parabolas, and Tangents}

\section{Line and ellipse}

The note continues the conic-section exercises.  A line and an ellipse can have three possible positional relationships:
\[
\begin{aligned}
  \Delta>0&:\ \text{two intersections},\\
  \Delta=0&:\ \text{tangent},\\
  \Delta<0&:\ \text{no real intersection}.
\end{aligned}
\]
This is obtained by substituting the line equation into the conic equation and examining the discriminant of the resulting quadratic.

For example, consider
\[
  y=2x-2,\qquad
  \frac{x^2}{5}+\frac{y^2}{4}=1.
\]
Substitution gives
\[
  \frac{x^2}{5}+\frac{(2x-2)^2}{4}=1,
\]
so
\[
  \frac65x^2-2x=0.
\]
Thus
\[
  x=0,\qquad x=\frac53,
\]
and the intersection points are
\[
  (0,-2),\qquad \left(\frac53,\frac43\right).
\]

\section{Tangents from a point to an ellipse}

Let
\[
  C:\frac{x^2}{a^2}+\frac{y^2}{b^2}=1
\]
and let \(P=(0,2)\).  A line through \(P\) has equation
\[
  y=kx+2.
\]
Substituting into the ellipse gives a quadratic in \(x\).  Tangency occurs exactly when the discriminant is zero.

For the example on the page, this produces
\[
  k=\pm\frac{\sqrt6}{3},
\]
so the two tangents are
\[
  y=\frac{\sqrt6}{3}x+2,\qquad
  y=-\frac{\sqrt6}{3}x+2.
\]

\section{Line and hyperbola}

For the hyperbola
\[
  \frac{x^2}{2}-\frac{y^2}{2}=1,
\]
and the line
\[
  y=x-1,
\]
substitution gives
\[
  x^2-(x-1)^2=2.
\]
Hence
\[
  2x-1=2,\qquad x=\frac32,\qquad y=\frac12.
\]
The line is not tangent; it meets the hyperbola at one point because the equation degenerates to a linear equation after substitution.  This is a useful warning: for hyperbolas, ``one intersection'' need not mean tangency, because the line may be parallel to an asymptotic direction.

\section{Parabola tangent through a point}

For the parabola
\[
  x^2=\frac12 y,
\]
we have
\[
  y=2x^2.
\]
A line through \(P=(2,6)\) may be written
\[
  y=kx+6-2k.
\]
Tangency means that
\[
  2x^2=kx+6-2k
\]
has a repeated root.  Thus
\[
  2x^2-kx-6+2k=0,
\]
and the discriminant condition is
\[
  k^2-4\cdot2(-6+2k)=0.
\]
Solving gives
\[
  k=12,\qquad k=-4.
\]
The tangent lines are therefore
\[
  y=12x-18,\qquad y=-4x+14.
\]

\section{A parabola and a circle}

The second page considers the parabola
\[
  C:x^2=2py,\qquad p>0,
\]
with focus \(F\), and a circle
\[
  M:x^2+(y+4)^2=1.
\]
The problem asks for \(p\) when the shortest distance between \(F\) and the circle is \(4\), and then asks for a maximal triangle area using two tangents from a point \(P\) on the circle.

Since the focus of \(x^2=2py\) is
\[
  F=\left(0,\frac p2\right),
\]
and the circle has center
\[
  O=(0,-4)
\]
and radius \(1\), the distance from \(F\) to the circle is
\[
  FO-1=\left(\frac p2+4\right)-1=\frac p2+3.
\]
Setting this equal to \(4\) gives
\[
  p=2.
\]
Hence the parabola is
\[
  x^2=4y.
\]

For a tangent point \(A=(x_1,x_1^2/4)\) on \(x^2=4y\), the tangent has equation
\[
  xx_1=2(y+x_1^2/4),
\]
or
\[
  y=\frac{x_1}{2}x-
rac{x_1^2}{4}.
\]
Similarly for a second tangent point \(B=(x_2,x_2^2/4)\).  The two tangents meet at
\[
  P=\left(\frac{x_1+x_2}{2},\frac{x_1x_2}{4}\right).
\]
Using the condition that \(P\) lies on the circle gives an optimization problem for the area of triangle \(PAB\).  The computation in the page reduces it to a one-variable expression and obtains the maximum
\[
  [PAB]_{\max}=20\sqrt5.
\]

\section{An ellipse with prescribed quadrilateral area}

The final exercise considers an ellipse
\[
  E:\frac{x^2}{a^2}+\frac{y^2}{b^2}=1,\qquad a>b>0,
\]
passing through \(A=(0,-2)\).  Four vertices form a quadrilateral of area \(4\sqrt5\).  Since \(A=(0,-2)\) lies on the ellipse, \(b=2\).  The area condition gives
\[
  2ab=4\sqrt5,
\]
so
\[
  a=\sqrt5.
\]
Therefore
\[
  E:\frac{x^2}{5}+\frac{y^2}{4}=1.
\]

For the line through \(P=(0,-3)\)
\[
  \ell:y=kx-3,
\]
intersecting the ellipse at \(B,C\), the equation is
\[
  \frac{x^2}{5}+\frac{(kx-3)^2}{4}=1.
\]
This gives a quadratic in \(x\).  The condition that the line meet the ellipse in two points is
\[
  k>1\quad\text{or}\quad k<-1.
\]
The problem then studies where the chords meet the line \(y=-3\) and imposes
\[
  |PM|+|PN|\le15.
\]
The computation on the page reduces this to
\[
  |k|\le3.
\]
Combining conditions gives
\[
  k\in[-3,-1)\cup(1,3].
\]

\section{The next layer}

The page is a concentrated exercise in the algebraic method for conics: substitute a line into a conic, read the discriminant, and interpret the result geometrically.  The only subtlety is that different conics behave differently: for ellipses, one intersection usually signals tangency; for hyperbolas, a line parallel to an asymptotic direction can produce a single finite intersection without being tangent.


\paragraph{Expanded synthesis.}
For this chapter, the elementary material should be read as a study of what remains invariant when coordinates change. The earlier sections (`Line and ellipse`, `Tangents from a point to an ellipse`, `Line and hyperbola`, `Parabola tangent through a point`, `A parabola and a circle`) compute with arrays, bases, pivots, determinants, or eigenvectors; the deeper object is a linear map together with the spaces it acts on. A matrix is a representation of that map after bases have been chosen. This is why formulas such as
\[
  A' = P^{-1}AP,\qquad \widetilde A = T^{-1}AS
\]
are not cosmetic: they distinguish an intrinsic morphism from a coordinate description.

The structural hierarchy is as follows. Kernels record what a map forgets, images record what it reaches, cokernels record obstructions to solvability, and quotients record information after an equivalence has been imposed. Determinants act on the top exterior power and therefore measure oriented volume; traces are invariant under cyclic rearrangement and become characters in representation theory; adjoints depend on an inner product and lead to self-adjoint spectral theory.

A useful way to return to the computations is to ask, for every formula, which invariant it is measuring. Row reduction measures rank and kernel; diagonalization measures decomposition into invariant subspaces; Gram--Schmidt measures orthogonal projection; positive definiteness measures ellipsoid geometry; tensor notation measures how components transform. The elementary methods are therefore the finite-dimensional laboratory in which duality, quotienting, functoriality, and spectral decomposition can be seen clearly.


\section{Discriminants as geometric detectors}

A line intersecting a conic becomes a quadratic equation after substitution.  The discriminant is therefore a geometric detector: positive means two intersections, zero often means tangency, and negative means no real intersection.  The word ``often'' matters because hyperbolas have asymptotic directions; a line can meet a hyperbola in only one finite point without being tangent.

This is why the same algebraic method must still be interpreted geometrically.  Solving equations gives candidates, but the conic type tells us what those candidates mean.
